\chapter{Known Limitations}

The development of the \eacsl plug-in is still ongoing. First, the whole \eacsl
reference manual~\cite{eacsl} is not yet supported. Which annotations are
already translated into \C code and which are not is defined in a separated
document~\cite{eacsl-implem}. Second, even if we do our best, bugs may exist. If
you find a new one, please report it on the bug tracking system (see Chapter 10
of the \framac User Manual~\cite{userman}). Third, there are some additional
known limitations, which could be annoying for the user in some cases, but are
hard to lift. Thus they could be there for a while. Lifting them could be part
of commercial supports\footnote{Contact us or read
  \url{http://frama-c.com/support.html} for additional details.}.

\section{Undefined Value}
\index{Undefined value}

\begin{itemize}
\item use of undefined values are not tracked in annotations
\item missing guards for \lstinline+\offset+, \lstinline+\block_length+ and
  \lstinline+texbase_addr+ (\lstinline+offset(p)+ must ensures that $p$ is
  valid)
\end{itemize}

\section{Incomplete Programs}

Section~\ref{sec:incomplete} explains how the \eacsl plug-in is able to handle
incomplete programs, which are either programs without main, or programs
containing functions without bodies.

However, if such programs contain memory-related annotations, the generated code
may be incorrect. That is made explicit by a warning displayed when the \eacsl
plug-in is running (see examples of Sections~\ref{sec:no-main} and
\ref{sec:no-code}).

\subsection{Program without Main}
\index{Program!Without main}
\label{sec:limits:no-main}

The generation is incorrect for every program without main containing a
memory-related annotations, except if the option
\optionuse{-}{e-acsl-full-mmodel} is provided.

Consider the following example.

\listingname{valid\_no\_main.c}
\cinput{examples/valid_no_main.c}

You can generate the instrumented program as follows.
\begin{shell}
\$ frama-c -e-acsl-full-mmodel -machdep x86_64 -e-acsl valid_no_main.c \
          -then-on e-acsl -print -ocode monitored_valid_no_main.i
<skip preprocessing commands>
[e-acsl] beginning translation.
<skip warnings about annotations from the Frama-C libc 
      which cannot be translated>
[kernel] warning: no entry point specified:
                  you must call function `e_acsl_global_init' and `__clean' by yourself.
[e-acsl] translation done in project "e-acsl".
\end{shell}

The last warning states an important point: if this program is linked against
another file containing a \texttt{main} function, then this main function must
be modified to insert a call to the function
\texttt{e\_acsl\_global\_init}\codeidx{e\_acsl\_global\_init} at 
the very beginning and a call to the function
\texttt{\_\_clean}\codeidx{\_\_clean} at the very end like in the following
example.

\listingname{modified\_main.c}
\cinput{examples/modified_main.c}

Then just compile and run it as explained in Section~\ref{sec:memory}.

\begin{shell}
\$ DIR=`frama-c -print-share-path`/e-acsl
\$ MDIR=\$DIR/memory_model
\$ gcc -o valid_no_main \$DIR/e_acsl.c \$MDIR/e_acsl_bittree.c \
      \$MDIR/e_acsl_mmodel.c monitored_valid_no_main.i modified_main.c
\$ ./valid_no_main
Assertion failed at line 11 in function f.
The failing predicate is:
freed: \valid(x).
\end{shell}

\subsection{Function without Code}
\label{sec:limits:no-code}
\index{Function!Without code}

\section{Recursive Function}
\index{Function!Recursive}

\section{Variadic Function}
\index{Function!Variadic}

\begin{itemize}
\item Not yet duplicated
\end{itemize}

\section{Function Pointer}
\index{Function!Pointer}
